\section{Martingales}

\begin{definition}
	Let $(\Omega, \cF, P)$ be a probability space, a martingale is a
	sequence $X_1, X_2, \dots$
	of random variables associated with $\sigma$-algebras $\cF_1, \cF_2, \dots$
	that satisfy:
	\begin{enumerate}[label=(\roman*)]
		\item $X_i \in L^1$ and $X_i \in \cF_i$.
		\item $\cF_i \subset \cF_{i+1}$.
		\item $\Ex[X_{i+1} | \cF_i] = X_i$.
	\end{enumerate}
\end{definition}

\begin{remark}
	If one changes the assumption of (iii) to be $\Ex[X_{i+1}|\cF] \leq
	X_i$ (resp. $\geq$), we call such sequence a supermartingale
	(resp. submartingale).
\end{remark}

\begin{theorem}(Doob's inequality)
	\label{trm:doob_ineq}
	Suppose $\{X_j\}$ is a martingale. Then
	\[
		\Pr(\sup_{j \leq n} |X_j| \geq l) \leq \frac{1}{l}
		\int\id[\sup_{j \leq n} |X_j| \geq l]|X_n|dP \leq \frac{1}{l}\Ex|X_n|.
	\]
\end{theorem}

\begin{proof}
	The proof is similar to Kolmogorov's maximal inequality. For $1
	\leq j \leq n$, let $E_j$ be the event
	$\{X_i < l, i < j\} \cap \{X_j \geq l\}$, that is the event that
	$X_j$ is the first to surpass $l$.
	Note that
	\begin{align}
		\Pr(\sup_{j \leq n} |X_j| \geq l) = \sum_{j \leq n} \Ex[\id[E_j]]
		\leq \sum_{j \leq n} \frac{1}{l}\Ex[\id[E_j]|X_j|] \leq \sum_{j
		\leq n} \frac{1}{l}\Ex[\id[E_j]|X_n|] = \Ex[\id[\sup_{j \leq n}
		|X_j| \geq l]|X_n|]
	\end{align}
	where the second inequality follows from the fact that $\Ex[X_n |
	\cF_j] = X_j$ and $|x|$ is a convex function.
\end{proof}

\begin{corollary}(Doob's p-inequality)
	\label{martingales:doob_p_ineq}
	Let $\{X_j\}$ be a martingale and $p > 1$. Then if, as before, $S_n
	= \sup_{j \leq n} |X_j|$, we have
	\[
		\Ex S_n^p \leq \bigg(\frac{p}{p-1}\bigg)^p\Ex |X_n|^p.
	\]
\end{corollary}
\begin{proof}
	Follows from the previous theorem and an integration by parts. Note
	that for any bounded positive r.v's $Y$ and $X$
	that satisfy $\Pr(Y \geq l) \leq \frac{1}{l}\Ex[\id[Y \geq l]X]$ we have
	\begin{align*}
		\int Y^p dP & = \int\int_0^Y py^{p-1} dy dP                                 \\
		& = p \int_0^\infty y^{p-1}\Pr(Y \geq y)dy                      \\
		& \leq p \int_0^\infty y^{p-1} \frac{1}{y}\Ex[\id[Y \geq y]X]dy \\
		& = p \int_0^\infty y^{p-2} \int \id[Y \geq y] X dP             \\
		& = p\int X \int_0^Y y^{p-2} dy = \frac{p}{p-1} \int X Y^{p-1}  \\
		& \leq \frac{p}{p-1}(\Ex X^p)^{1/p} (\Ex Y^p)^{\frac{p-1}{p}}
	\end{align*}
	So that $\Ex Y^p \leq (\frac{p}{p-1})^p \Ex X^p$. We can then
	consider a truncation argument to prove the result
	for unbounded positive random variables.
\end{proof}
\begin{remark}
	This corollary is very cool on its own, it says that if some $X_n$
	in the martingale belongs to $L^p$, then
	all previous r.v's belong as well.
\end{remark}

\begin{exercise}
	Consider the martingale of doubling one's bet. In the space $([0,1],
	\cB([0,1]), dm)$,
	we consider the random variables $X_n \in \sigma([j/2^n, (j+1)/2^n])$
	\[
		X_n(\omega) =
		\begin{cases}
			2^n & \text{ if } \omega \in [0,2^{-n}] \\
			0   & \text{ elsewhere}
		\end{cases}
	\]
	It is easy to see that they define a martingale, all in $L^1$,
	$\Ex{X_n} = 1$, but $\Ex \sup_{j \leq n} X_j = (n+1)/2$.
	So no alternative to [\ref{martingales:doob_p_ineq}] exists in $L^1$.
\end{exercise}

\begin{exercise}
	If $\{X_n\}$ is a martingale with all differences $Y_n = X_{n} -
	X_{n-1}$ in $L^2$, show that for $n \neq m$,
	the differences are orthogonal, that is $\Ex Y_nY_m = 0$. Therefore
	\[
		\Ex X_n^2 = \Ex X_0^2 + \sum_{j = 1}^n \Ex[Y_n^2].
	\]
	Conclude that if $\sup \mathrm{Var}X_n < \infty$, $X_n$ converges in
	$L^2$ to a r.v $X$.
\end{exercise}

\begin{proof}
	It follows from the properties of condition expectation. Suppose $n
	> m$, note that
	\\
	\begin{align*}
		\Ex[(X_n - X_{n-1})(X_m - X_{m-1})] & = \Ex\Ex[(X_n - X_{n-1})(X_m -
		X_{m-1}) | \cF_{n-1}]
		\\
		& = \Ex[\Ex[X_n - X_{n-1}| \cF_{n-1}](X_m - X_{m-1})] \\
		& = \Ex[(X_{n-1} - X_{n-1})(X_m - X_{m-1})] = 0
	\end{align*}
	So the first result follows from that and the cancellation of cross
	terms. The conclusion holds as $\sum_n \Ex Y_n^2$ is bounded,
	so $\{X_n\}$ is a cauchy seq. in $L^2$.
\end{proof}

From now on, we denote as $\cF$ the smallest $\sigma$-field
containing $\{\cF_n\}$. If $X \in \cF$, then $X_n = \Ex[X | \cF_n]$
is a (Doob) martingale,
and, by Jensen's ineq, if for some $p \geq 1$ $X \in L^p$ then $X_n
\in L^p$ for all $n$.

\begin{theorem}
	If $p \geq 1$ and $X \in L^p$, then $\lim_{n\to\infty} ||X - X_n||_p = 0$.
\end{theorem}

\begin{proof}
	By the last exercise, the result is true if $X$ is bounded. If not,
	approximate $X$ by a bounded $X'$
	and note that if $||X - X'|| < \eps$, then $||\Ex[X_n | \cF_n]  -
	\Ex[X'|\cF_n]||_p \leq \eps$.
	So we can do a two step approximation.
\end{proof}

\begin{theorem}
	If, for some $p > 1$, a martingale $\{X_n\}$ is bounded in $L^p$,
	that is $\sup_n ||X_n||_p < \infty$,then there exists $X \in L^p$ s.t.
	$X_n = \Ex[X | \cF_n]$ for all $n$. In particular $||X_n - X||_p \to 0$.
\end{theorem}

\begin{proof}
	As $p > 1$, we can choose a subseq $X_{n_j}$ that weakly converges
	to a limit $X$ in $L^p$ (just like Egorov's).
	By noting that for any $A \in \cF_m$ we have $\id[A] \in L^q$, we see
	\[
		\Ex[\id[A]X_m] = \lim_{j \to \infty} \Ex[\id[A]X_{n_j}] = \Ex[\id[A]X].
	\]
	So $X_m = \Ex[X | \cF_m]$.
\end{proof}

\begin{theorem}
	Let $X \in L^p$ for $p \geq 1$. Then the martingale $\{X_n\}$ with
	$X_n = \Ex[X|\cF_n]$ converges to $X$ for
	almost all $\omega$ with respect to $P$.
\end{theorem}
\begin{proof}
	It suffices to prove for $p = 1$. This proof is a bit nuts, we
	consider the set $\cM \subset L^1$
	of random variables $Y$ s.t $Y_n \to Y$ a.e. We show that this set
	is dense and closed.
	It is easy to see that it is dense as all variables belonging to
	$\cF_n$ are in it.
	To see how the set is closed, pick a sequence $Y_j \to X$ in $L^p$.
	Note that, as $X \in L^1$, by [\ref{martingales:doob_p_ineq}], $X_n$
	is an almost surely bounded
	sequence so it suffices to show $\Pr(\limsup{X_n} - \liminf{X_n} >
	\eps) = 0$ a.e. Set $Y_{j,n} = \Ex[Y_j | \cF_n]$,
	and note that, for every $j$
	\[
		\limsup_n{X_n} - \liminf_n{X_n} \leq \limsup_n{Y_{j,n}} -
		\liminf_n{Y_{j,n}} + 2\sup_n{|X_n - Y_{j,n}|} = 2\sup_n{|X_n - Y_{j,n}|}.
	\]
	So
	\[
		\Pr(\limsup_n{X_n} - \liminf_n{X_n} > \eps) \leq \Pr(\sup_n |X_n -
		Y_{j,n}| > \eps) \leq \Pr(|X - Y_j| > \eps) \leq
		\frac{1}{\eps}\Ex[|X - Y_j|] \to 0
	\]
	when $j$ tends to infinity.
\end{proof}

\begin{theorem}

\end{theorem}
